% main.tex - 科研项目学习笔记模板
% !TeX root = main.tex
%%%%%%%%%%%%%%%%%%%%%%%%%%%%%% DOCUMENT
\documentclass[11pt]{article}

%%%%%%%%%%%%%%%%%%%%%%%%%%%%%% PACKAGES

% 中文支持（XeLaTeX 编译）
% \usepackage[UTF8]{ctex}
% \usepackage{xeCJKfntef}

% \setCJKmainfont{HanziPen SC}
% \setmainfont{HanziPen SC}

% 页面设置
\usepackage[a4paper, left=25mm, right=25mm, top=25mm, bottom=25mm]{geometry}

\PassOptionsToPackage{dvipsnames,svgnames,x11names}{xcolor}
\usepackage{xcolor}

% 数学环境及符号
\usepackage{amsmath, amssymb, amsfonts, amsthm,amsopn}
\usepackage{tensor}              % 张量指标管理
\usepackage{mathtools}           % amsmath增强
\usepackage{physics}             % 物理公式快捷命令
\usepackage{bbold}               % 数学黑体
\usepackage{dsfont}              % 另一种数字体
\usepackage[mathscr]{eucal}     % 花体字母
\usepackage{tensor}              % 张量指标管理
\usepackage{simpler-wick}       % Wick记号
\usepackage{mathrsfs}            % 另一种花体字母

% 颜色与图形相关
\usepackage{graphicx}           % 插图支持
\usepackage{float}              % 浮动体控制
\usepackage{tikz}               % 绘图库
\usetikzlibrary{math}           % tikz数学扩展
\usepackage{geometry}
% 表格与列表
\usepackage{makecell}           % 表格多行换行
\usepackage{multicol}           % 多栏排版
\usepackage{colortbl}           % 表格颜色
\usepackage{enumitem}           % 列表自定义

% 其他辅助
\usepackage{framed}             % 有边框环境
\usepackage{tcolorbox}          % 灵活盒子环境
\tcbuselibrary{breakable}       % 盒子内容分页
\usepackage{thmtools}           % 定理环境管理
\usepackage{thm-restate}        % 定理重述
% \usepackage{showlabels}         % 显示标签，调试用（完成后可注释）
\usepackage[normalem]{ulem}     % 下划线、删除线
\usepackage{hyperref}           % 超链接（最后加载）
\usepackage{cleveref}           % 智能引用（紧跟hyperref）
\usepackage{soul}
\definecolor{lightred}{rgb}{1,.8,.8}
\sethlcolor{lightred}
\usepackage{hyperref}

\usepackage{natbib}
\bibliographystyle{plain}

% 自定义宏包
\usepackage{macros}

% 一个中文可以高亮的包
\usepackage{cjkhl}
\definecolor{lightblue}{rgb}{.8,.8,1}

%%%%%%%%%%%%%%%%%%%%%%%%%%%%%% 自定义命令
\newcommand{\tml}{Teichmüller space}
\newcommand{\hil}{Hilbert space}
\newcommand{\mtc}{Modular Tensor Category}

%%%%%%%%%%%%%%%%%%%%%%%%%%%%%% BEGINNING OF THE DOCUMENT

\begin{document}

\title{\boldmath Bosonization of Majorana CFT with Boundaries}
\author{Yu Liu}

\maketitle

\begin{abstract}
  This is my own note on the topic of Majorana CFT and Ising CFT with
  boundaries.

\end{abstract}

\tableofcontents

\newpage
\section{Preliminaries 1: Ising CFT and BCFT}\label{sec:Boundary
Ising CFT} % (fold)
\input{pre.tex}
% section Boundary Ising CFT (end)

\section{Preliminaries 2: Majorana CFT}\label{sec:Majorana CFT} % (fold)
\input{majorana_cft.tex}

% section Majorana CFT (end)

\newpage
\section{Majorana BCFT}\label{sec:Free Majorana Fermion BCFT} % (fold)
\input{majorana.tex}
% section Free Majorana Fermion BCFT (end)

\newpage
\section{Bosonization of Majorana BCFT}\label{sec:Bosonization of
Majorana BCFT} % (fold)
\input{bosonization.tex}
% section Bosonization of Majorana BCFT (end)

% \newpage
% \section{Bosoniztaion from Exact Partition
% Function}\label{sec:Bosoniztaion from Exact Partition Function} % (fold)
% \input{bosonization_exact.tex}
% % section Bosoniztaion from Exact Partition Function (end)

\newpage
\bibliography{references}

\end{document}
